\section{Comparison of the two higher-genus actions}
The explicit action on the groups with zero or one marked point agrees,
as an outer action, with the action constructed from quilts. The following
calculation gives the change of representative.

\begin{proposition}\label{prop:ns-lns-comparison}
Use the compatible standard markings of Nakamura--Schneps
\cite[Theorem 11.2]{NS2000} and Lochak--Nakamura--Schneps
\cite[Theorem 1.1]{LNS}. For $F=(\lambda,f)\in\widehat{\operatorname{GT}}$,
write $\rho=\rho_2(F)$ and
\[
 w_1=1,\qquad w_j=(a_1\cdots a_{j-1})^j\quad(2\leq j\leq2g),
 \qquad
 K_F=\prod_{j=2}^{2g}w_j^{\,(-1)^{j+1}4\rho}.
\]
Let $A_F^{\mathrm{NS2000}}$ denote the standard quilt representative and
$A_F^{\mathrm{LNS}}$ the representative given by the explicit chain
formulas. On each of $\widehat\Gamma_{g,0}$ and
$\widehat\Gamma_{g,1}$ one has
\[
 A_F^{\mathrm{NS2000}}=\operatorname{Ad}(K_F)\circ A_F^{\mathrm{LNS}}.
\]
In particular the two outer actions are equal. If $\rho_2(F)=0$, these
standard representatives themselves are equal.
\end{proposition}
\begin{proof}
The equality of the Nakamura--Schneps subgroup with
$\widehat{\operatorname{GT}}$ is used as established background
\cite[Corollary 21]{H17}; it ensures that their formulas apply to every
$F$ in the statement.

The elements $a_1,\ldots,a_{2g}$ satisfy the braid relations. The full
braid twist $w_j$ is central in the subgroup generated by
$a_1,\ldots,a_{j-1}$. Hence the $w_j$ commute with one another, and
$w_j$ commutes with $a_i$ whenever $j\ne i$. The usual full-twist
factorization gives
\[
 w_i=y_iw_{i-1},\qquad
 y_i=a_{i-1}\cdots a_1^2\cdots a_{i-1}.
\]
The element $w_{i-1}$ commutes both with $a_i$ and with $y_i$.
If $z$ centralizes $a,b$ and $f$ belongs to the closed commutator
subgroup of a free profinite group on two generators, then
$f(a,bz)=f(a,b)$. Indeed the homomorphism from
$\langle a,b\rangle\times\langle z\rangle$ given by multiplication
sends the second coordinate of a commutator word to the identity;
the identity extends from commutator words by continuity.
Consequently
\[
 f(a_i^2,w_i)=f(a_i^2,y_i),\qquad
 f(w_i,a_i^2)=f(y_i,a_i^2).
\]

Set $c_i=f(a_i^2,y_i)$. All $w_j$ with $j\ne i$ commute with
$c_i$ and with $a_i$. The explicit chain action is
$A_F^{\mathrm{LNS}}(a_i)=c_i a_i^\lambda c_i^{-1}$.
Conjugating by $K_F$ therefore gives
\[
 w_i^{\,(-1)^{i+1}4\rho}
 f(a_i^2,w_i)a_i^\lambda f(w_i,a_i^2)
 w_i^{-(-1)^{i+1}4\rho},
\]
which is the standard quilt formula. For $a_1$ both maps give
$a_1^\lambda$, and $K_F$ commutes with $a_1$.

For $g\geq2$, the remaining generator is $d=d_2$, a curve of the
standard pants decomposition. In the same marking,
\[
 w_2=a_1^2,\qquad
 w_{2i}=d_i d_{-i}\quad(i\geq2),\qquad
 w_{2i-1}^{2}=e_{i-1}\quad(i\geq2).
\]
All curves appearing on the right belong to that pants decomposition.
Since $4\rho$ is divisible by two, these identities express $K_F$ as
a product of powers of its mutually commuting pants twists. Thus
$K_F$ commutes with $d$. Both representatives take $d$ to $d^\lambda$.
Equality on the displayed generating set proves the assertion; in
genus one the chain generators already suffice. The same identities
survive the relation filling the marked point, so the argument also
applies to $\widehat\Gamma_{g,0}$.
\end{proof}

This is an equality on the overlapping families of full profinite
mapping class groups. A comparison with another completed-tower
construction additionally requires that construction's action formulas
or a stated uniqueness theorem for its objects and morphisms. The
existence of actions with the same genus-zero restriction alone does
not identify their chosen representatives on every object of a tower.

\section{The scope of inertia-preserving identifications}
For a generic surface $S$ of type $(g,n)$, namely
$g>0$, $3g-3+n>2$, and $(g,n)\ne(2,0)$, we use the established
identification
\[
 \operatorname{Out}^{*}(\check\Gamma(S))
       \cong\widehat{\operatorname{GT}}.
\]
Here the star requires preservation of the twist inertia subgroups.
This is \cite[Corollary 22]{H17}, using Lochak's compatible-pair
classification \cite[Definition 5.4, Proposition--Definition 5.6,
and Theorem 5.10(ii)]{Lochak2025}. It is a cited result about the
procongruence completion. Passing from nonseparating inertia to all
inertia is a further assertion.

Boggi states this further assertion as
$\operatorname{Aut}^{I_0}(\check\Gamma(S))=
\operatorname{Aut}^{I}(\check\Gamma(S))$ in
\cite[Theorem 3.12 and Corollary 3.15]{Boggi}.
The proof proceeds through the nonseparating curve complex and requires
recognition of the multicurves with connected complement, as used in
\cite[Remark 3.11]{Boggi}. Commutation of individual twists recognizes
disjointness. It does not by itself recognize connectedness of the
complement: a bounding pair and a pair with connected complement both
consist of disjoint nonseparating curves. Thus the latter recognition
step requires its own justification.

\begin{proposition}\label{prop:nonunique-stabilizer-extension}
Let $\sigma=\{\gamma_1,\ldots,\gamma_{g-1}\}$ be a cut system with
connected complement, and let $H$ be its oriented stabilizer in the
full profinite pure mapping class group. Let $H_i$ be the oriented
stabilizer of $\sigma\setminus\{\gamma_i\}$. The identity
of $H$ has more than one inertia-preserving extension to $H_i$.
\end{proposition}
\begin{proof}
The twist $\tau_{\gamma_i}$ is central in $H$, since every element
of the oriented stabilizer fixes its twist. Thus the identity and
$\operatorname{Ad}(\tau_{\gamma_i})$ restrict to the same automorphism
of $H$. Both are inertia-preserving automorphisms of $H_i$.
Choose a simple curve $\beta_i$ meeting $\gamma_i$ once and disjoint
from the other curves of the cut system. Its twist belongs to $H_i$.
On the symplectic plane generated by these two curves, their twist
matrices are
\[
 U=\begin{pmatrix}1&1\\0&1\end{pmatrix},\qquad
 V=\begin{pmatrix}1&0\\-1&1\end{pmatrix}.
\]
They do not commute, already after reduction modulo three. Thus
$\tau_{\gamma_i}$ is not central in $H_i$, and its inner automorphism
is not the identity. This finite symplectic quotient also shows that
the distinction survives profinite completion.
\end{proof}

Accordingly, uniqueness as a genuine automorphism in
\cite[Lemma 6.16]{Boggi} fails as written. This counterexample does
not refute a statement of uniqueness up to inner automorphisms.
There is a separate group identification issue in the proof of
\cite[Proposition 6.9]{Boggi}: the quotient
$B_{2g+1}/\langle z^2\rangle$ obtained by filling a distinguished
Weierstrass boundary point preserves that point. It therefore gives
the corresponding stabilizer in the closed hyperelliptic group,
whose index is $2g+2$, rather than the whole closed hyperelliptic
group. Neither observation is a counterexample to the earlier
nonseparating-inertia theorem; they concern distinct proof assertions.
